% QMB_n12_grenze.tex — Kapitel 12: Grenzen des Quantenrechnens
% Inhaltliche Grundlage: Dok. 075, 076, 147, 316, 343; eigene Darstellung
\chapter{Grenzen des Quantenrechnens: RSA, Weyl und der Auflösungsboden}
\label{ch:grenze}

\section{Was Quantencomputer leisten und was nicht}

Dieses Kapitel zieht drei Grenzen, die der Korpus mit Begründung
festlegt. Sie sind keine spekulativen Einschränkungen, sondern
Konsequenzen mathematischer Sätze oder empirischer Befunde.

\section{Grenze 1: RSA-Faktorisierung und der arithmetische Träger}

RSA (Rivest, Shamir, Adleman 1977) beruht auf der Schwierigkeit der
Faktorisierung großer Semiprimzahlen $N = p\cdot q$. Shors Algorithmus
löst das Problem in polynomialer Zeit --- in Prinzip. Die Frage ist, was
der Geometrieparameter $\xi$ dabei leisten kann.

Das Ergebnis des Korpus ist klar: \textbf{$\xi$ hat in $\mathbb{Z}/N\mathbb{Z}$
keine Funktion.}\status{K} Die Periode $r$ der modularen Potenzfunktion
$f(x) = a^x\bmod N$ ist durch $r = \mathrm{lcm}(p-1,q-1)$ bestimmt ---
eine individuelle arithmetische Eigenschaft von $N$, nicht ein
geometrischer Parameter. Shors Algorithmus nutzt diese Periodenstruktur;
sie ist eine Eigenschaft der Primzahlen, nicht der FFGFT-Geometrie.

Die Tests in Dok.~075 und Dok.~076 verwenden Rasterparameter $\sigma$
mit Werten wie $1/42$, $1/100$, $1/1000$ für Periodensuchverfahren.
Diese Größen sind numerische Abtastweiten, keine physikalischen
Parameter; sie dürfen nicht mit $\xi = 4/30000$ gleichgesetzt werden.
Das Resonanzbild ist prinzipiell richtig --- Periodenfindung ist ein
spektrales Problem ---, aber die Generatoren der Periode sind die
Primzahlen, nicht ein einstellbarer Geometrieparameter.

Ein Vorteil gegenüber klassischen Algorithmen muss auf der
Shor-Faktorisierungskurve sichtbar werden: Der Fehlerkorrektur-
Schwellenwert muss für relevante Schlüssellängen (2048 Bit) gehalten
werden; die Aufbereitung muss in Superposition ausgeführt werden; ein
RSA-relevantes $N$ muss erfolgreich faktorisiert werden. Nichts davon
ist Stand 2026 realisiert.\status{X}

\section{Grenze 2: Weyl-Obstruktion und Zähldichte}

Der Weyl-Satz bestimmt, wie viele Eigenwerte ein kompakter Zustandsraum
unterhalb einer Schranke $\lambda$ haben kann:
$N(\lambda)\sim C\cdot\lambda^{d/2}$
--- ein Potenzgesetz in der Dimension $d$.\status{K}

Die Primzahlfunktion $\pi(T)$ und die Riemannschen Nullstellen wachsen
dagegen wie $T\ln T$ --- keine Potenz, sondern quasi-logarithmisch. Das
bedeutet: Ein endlicher kompakter Zustandsraum (ein $n$-Qubit-Register
ist ein solcher) kann ein solches Spektrum nicht vollständig tragen.
Die Weyl-Obstruktion schließt drei geometrische Zugänge zur
Riemann-Hypothese aus:\status{K} keine kompakte Geometrie kann
$T\ln T$-Wachstum aus einem Laplace-Spektrum liefern; $\xi$ verschiebt
die kritische Linie nicht; die $\xi$-Leiter diskretisiert die
Nullstellen nicht.

Shors Schaltkreismodell ist digital und von der Weyl-Grenze direkt nicht
betroffen. Aber ein physikalischer Quantencomputer ist ein Resonanzsystem;
solange das Schwellentheorem für relevante Skalierungen nicht demonstriert
ist, liegt er innerhalb der Resonanzgrenzen.\status{S}

\section{Grenze 3: Der Auflösungsboden}

Die $\xi$-Korrektur zu Bell-Korrelationen liegt bei $\Delta\mathrm{CHSH}
\approx\xi/(2\pi)\approx2\times10^{-5}$ pro Messung.\status{S} Das ist
der Auflösungsboden: Unterhalb dieser Schwelle ist keine FFGFT-spezifische
Abweichung von der Standard-QM messbar. Alle NISQ-Prozessoren (NISQ =
Noisy Intermediate-Scale Quantum) operieren mit physikalischen Fehlerraten
$\sim10^{-2}$, drei Größenordnungen über dem $\xi$-Signal.

Das bedeutet für den Quantenrechner: Die $\xi$-Korrekturen zu Algorithmen
--- Bell-Dämpfung in Shor, modifizierte Gatterwinkel durch $K_{\mathrm{frak}}$
--- sind unter NISQ-Rauschen begraben.\status{K} Weder eine Verlangsamung
noch eine Beschleunigung von Quantenalgorithmen durch $\xi$ ist mit
aktueller Hardware messbar. Das ist kein Misserfolg, sondern eine
Präzisionsaussage.

Zukünftige Hochpräzisions-Bell-Tests (schlupflochfrei, Detektions-
effizienz $>99\,\%$, optische Plattform) könnten die Grenze erreichen.
Das photonische System (Kapitel~\ref{ch:hardware}) ist der natürliche
Kandidat.

\section{Was bleibt}

Von den vier Parallelstrukturen zwischen FFGFT und Quantenalgorithmen
bleiben nach dieser Grenzziehung:

Die Übersetzbarkeit (Kapitel~\ref{ch:hilbert}) ist vollständig und
symmetrisch. Jede QM-Rechnung ist in FFGFT formulierbar; jede FFGFT-Aussage
über Quantensysteme ist in der QM abbildbar. Das ist keine Einschränkung,
sondern eine Stärke: Praktiker können ihren vertrauten Formalismus
behalten.

Der Mehrwert der FFGFT liegt nicht in anderen algorithmischen Ergebnissen,
sondern in der Herkunft der Eingabewerte: Massen, Kopplungen und
$\alpha$ aus derselben Geometrie, die den Quantenformalismus trägt. Das
ist die strukturelle Asymmetrie, die Kapitel~\ref{ch:hilbert} festgehalten
hat.

\quelle{Dok.~075 (RSA-Verfahren), Dok.~076 (RSA-Periodentests), Dok.~147 §8 (IBM-Validierung), Dok.~316 (Riemann, Weyl-Grenze), Dok.~343 (Auflösungsboden)}
